\section{Resultados de la simulación}\label{sec:resultados}

\subsection{Variables de canal, identificación y especificación}\label{sec:res-principal}

En los ocho escenarios de la tabla~\ref{tab:principal}, las medias de los estimadores quedan cerca de su límite predicho (todos los $|z|\le1,3$ salvo con solapamiento medio, que se comenta a continuación). Con muestras pequeñas fallan la dispersión y la cobertura, no la media (sección~\ref{sec:res-bootstrap}). Una discrepancia en el solapamiento medio ($z=-3,0$ en el Hájek y $-2,1$ en el DR) no reapareció en una réplica con 200 semillas nuevas ($z=+0,8$ y $+1,5$; con las 400 réplicas, $-1,4$ y $-0,3$). Esa réplica no estaba planificada de antemano; con unas 40 comparaciones, la probabilidad de algún $|z|\ge3$ por azar ronda el 10\,\%.

\begin{table}[ht]
\centering\small
\caption{Resultados principales. Tasas en \%; entre paréntesis, $z=$ (media $-$ predicción)/(error de Monte Carlo). DE: desviación típica del DR entre réplicas. Cobertura del intervalo de Wald del DR al 95\,\%, respecto de $\theta$. En los escenarios con $C$, el $\neff/n_0$ teórico se calcula con la parte con soporte; sale prácticamente igual que sin $C$ (0,217), porque $\PP(H=0\mid V,D=1)$ varía poco en la región objetivo.}\label{tab:principal}
\resizebox{\textwidth}{!}{%
\begin{tabular}{@{}lcccccc c@{}}
\toprule
Escenario & $\theta$ & Predicción & Hájek & DR & DE & Cobertura & $\neff/n_0$ (teórico) \\
\midrule
Base & 0,958 & 0,958 & 0,958 (0,0) & 0,961 (0,4) & 0,112 & 95,5\,\% & 0,14 (0,22) \\
Con $C$ & 0,958 & H 0,940; DR 0,958 & 0,939 ($-$0,1) & 0,958 (0,0) & 0,107 & 96,5\,\% & 0,12 (0,22) \\
$C$ confunde, sin $C$ & 1,712 & 1,314 & 1,313 ($-$0,1) & 1,320 (0,6) & 0,132 & 15,5\,\% & 0,14 (0,22) \\
$C$ confunde, con $C$ & 1,712 & H 1,281; DR 1,314 & 1,280 (0,0) & 1,320 (0,6) & 0,129 & 17,5\,\% & 0,12 (0,22) \\
Solapamiento medio & 1,178 & 1,178 & 1,151 ($-$3,0) & 1,158 ($-$2,1) & 0,136 & 94,5\,\% & 0,10 (0,14) \\
Solapamiento bajo & 1,274 & 1,274 & 1,283 (0,5) & 1,298 (1,1) & 0,300 & 91,0\,\% & 0,03 (0,03) \\
Confusión no observada & 2,697 & 1,193 & 1,198 (0,6) & 1,204 (1,3) & 0,127 & 0,0\,\% & 0,14 (0,20) \\
3\,000 clientes & 0,958 & 0,958 & 0,887 ($-$1,2) & 0,883 ($-$1,2) & 0,912 & 70,0\,\% & 0,02 (0,22) \\
\bottomrule
\end{tabular}}
\end{table}

\begin{itemize}
  \item \emph{Variable de canal.} Con $C$, la masa sin soporte estimada fue $\hat\varepsilon=0,326$, frente a $\kappa=0,335$ (Proposición~\ref{prop:rompe}), y el Hájek converge a $\thac$ (el límite de la ponderación de la sección~\ref{sec:canal} coincide aquí con $\thac$ porque $\delta=0$), con el sesgo predicho de $-0,018$ puntos. La predicción original asignaba también al DR el límite $\thac$ y falló ($z=2,4$ y $4,3$). El DR extrapola el modelo de resultado a la parte sin soporte (Corolario~\ref{cor:drsinsoporte}); la predicción derivada de ello (0,958 y 1,314) se añadió tras el escenario «Con $C$» y se registró antes de ejecutar «$C$ confunde, con $C$», donde se cumplió ($z=0,6$).
  \item \emph{Confusión por el historial.} Sin $C$, el sesgo es el predicho: $-0,40$ puntos, un 23\,\% de $\theta$. Con $C$ tampoco se recupera $\theta$, como dice la sección~\ref{sec:canal}.
  \item \emph{Identificación.} Sin confusión ni variable de canal, Hájek y DR recuperan $\theta$ y el intervalo cubre el 95,5\,\%.
  \item \emph{Confusión no observada.} Los estimadores convergen al límite predicho por los datos observables, a menos de la mitad de $\theta$, con cobertura nula. El error típico está bien calibrado (97\,\% de la dispersión real) y los diagnósticos se parecen a los del escenario base: nada en los datos delata el problema.
  \item \emph{Solapamiento.} Con solapamiento bajo ($\chi^2\approx36$, sección~\ref{sec:diseno}), la desviación típica casi se triplica, el error típico estimado se queda en el 87\,\% de la dispersión real y la cobertura baja al 91\,\%.
\end{itemize}

\begin{table}[ht]
\centering\small
\caption{Especificación (escenario base). Tasas en \%.}\label{tab:especificacion}
\begin{tabular}{@{}llllcc@{}}
\toprule
Clasificador & Modelo de resultado & Estimador & Predicción & Media ($z$) & Cobertura \\
\midrule
LightGBM & LightGBM & DR & 0,958 & 0,961 (0,4) & 95,5\,\% \\
LightGBM & lineal en $V_1$ & DR & 0,958 & 0,955 ($-$0,3) & 95,0\,\% \\
Logística lineal & LightGBM & DR & 0,958 & 0,946 ($-$2,3) & 88,0\,\% \\
Logística lineal & --- & Hájek & 1,024 & 1,023 ($-$0,4) & --- \\
Logística lineal & lineal en $V_1$ & DR & 1,157 & 1,150 ($-$1,5) & 13,5\,\% \\
--- & lineal en $V_1$ & Regresión & 1,069 & 1,075 (1,4) & --- \\
\bottomrule
\end{tabular}
\end{table}

Con el clasificador flexible y el modelo de resultado mal especificado (lineal), el DR no tiene sesgo y cubre el 95\,\%; con los dos mal, converge al límite predicho, lejos de $\theta$ (tabla~\ref{tab:especificacion}). Con el clasificador mal especificado (lineal) y el modelo de resultado flexible, queda un sesgo de $-0,012$ puntos y una cobertura del 88\,\%: el modelo de resultado es un LightGBM, no $m_0$, y el producto de errores no es despreciable a estos tamaños. Aun así, el DR reduce en un 82\,\% el sesgo del Hájek correspondiente.

\subsection{Recorte}\label{sec:res-recorte}

Con los pesos verdaderos y $m_0$, recortar en los percentiles 99, 95 y 90 baja la estimación en $0,37$, $3,13$ y $7,81$ milésimas de punto porcentual, como predice la Proposición~\ref{prop:recorte} ($0,37$, $3,18$ y $7,98$): los perfiles con más peso responden más. Con pesos estimados el signo se invierte ($+10,9\pm0,7$ en el percentil 90, positivo en las 60 réplicas). Dentro del DR, el efecto del recorte no se distingue de cero y reduce la raíz del error cuadrático medio (de 0,300 a 0,163 puntos con solapamiento bajo). Su error típico debe calcularse con los pesos sin recortar: con los recortados, la cobertura cae al 78,5\,\% con solapamiento bajo, y aun con los sin recortar es solo del 78--88,5\,\% con 3\,000 clientes. La inversión se analiza en la nota complementaria.

\subsection{Pesos con corte conocido}\label{sec:res-corte}

Comparamos el estimador con los pesos de la Proposición~\ref{prop:corte} con el estándar en las mismas réplicas: mismos datos, mismo LightGBM y mismo modelo de resultado; solo cambian los pesos. La propensión de canal se estima añadiendo las llamadas premium de fuera del grupo objetivo, que el estimador estándar no usa. Las predicciones se registraron antes de ejecutar estos escenarios.

\begin{table}[ht]
\centering\small
\caption{Pesos con corte conocido en el escenario base, por tamaño. 200 réplicas; $\theta=0,958\,\%$. DE en puntos. La última columna es el efecto del recorte del Hájek en el percentil 90 calculado con $m_0$, en milésimas, sobre estas 200 réplicas (el $+10,9$ del texto procede de otro conjunto de 60); con los pesos verdaderos es $-7,8$.}\label{tab:corte}
\begin{tabular}{@{}llccccc@{}}
\toprule
Clientes & Pesos & DR (\%) & DE & Cobertura & $\neff/n_0$ & Recorte p90 \\
\midrule
30\,000 & estándar & 0,961 & 0,112 & 95,5\,\% & 0,14 & $+10,4$ \\
30\,000 & corte conocido & 0,960 & 0,098 & 92,5\,\% & 0,20 & $-6,2$ \\
10\,000 & estándar & 0,962 & 0,294 & 88,0\,\% & 0,07 & $+19,8$ \\
10\,000 & corte conocido & 0,957 & 0,195 & 91,0\,\% & 0,17 & $-4,8$ \\
3\,000 & estándar & 0,883 & 0,912 & 70,0\,\% & 0,02 & $+50,0$ \\
3\,000 & corte conocido & 0,944 & 0,371 & 87,0\,\% & 0,11 & $-10,2$ \\
\bottomrule
\end{tabular}
\end{table}

Se cumplen cinco de las seis predicciones (tabla~\ref{tab:corte}): no hay sesgo apreciable; $\neff/n_0$ se acerca al 0,22 teórico; la dispersión baja un 12, un 34 y un 59\,\%; el recorte recupera el signo de los pesos verdaderos; y la cobertura mejora con 10\,000 y 3\,000 clientes. Falla la cobertura con 30\,000: predijimos entre el 93 y el 97\,\% y sale $92,5\,\%\pm1,9$, porque el error típico se queda en el 93\,\% de la dispersión real. Con 10\,000 clientes, la desviación típica queda un 15\,\% por encima de la que predice el ritmo $\sqrt n$ a partir del mismo estimador con 30\,000 ($0,098\cdot\sqrt3=0,170$), frente al 52\,\% del estándar.

Un análisis posterior, no registrado de antemano, separa las dos novedades: imponer el corte y usar las unidades premium de fuera del grupo objetivo. Imponer el corte sobre el clasificador estándar, sin datos nuevos, no reduce la dispersión (0,117, 0,298 y 0,933 puntos) y corrige solo en parte el signo del recorte ($+1,0$, $+2,7$ y $+9,6$). La mejora viene sobre todo de estimar la propensión de canal con muchos más casos positivos; el corte conocido es lo que permite usarla.

En los demás escenarios (apéndice~\ref{app:extra}), los estimadores coinciden con los mismos límites que el estándar, y la ganancia de precisión se mantiene con buen solapamiento pero desaparece con solapamiento medio y bajo.

\subsection{Sensibilidad}\label{sec:res-sensibilidad}

El $\Lambda$ mínimo para el que el intervalo poblacional $[\theta_0/\Lambda,\Lambda\theta_0]$ contiene $\theta$ es $\theta/\theta_0$: 1,30 cuando el historial confunde y 2,26 con confusión no observada. La cobertura del intervalo combinado $[\hat\ell/\Lambda,\hat u\Lambda]$ sube deprisa alrededor de ese valor (tabla~\ref{tab:sensibilidad}); queda por debajo de la aproximación normal porque el error típico es menor en las réplicas con estimación baja (correlación de 0,64 y 0,70). La condición de resultado raro de la Proposición~\ref{prop:sensibilidad}, comprobada mediante la cota $\PP(Y>0\mid V,D=0)\le g(V)=\E[Y\mid V,D=0]$ en el percentil 99,99 de $g$ en el canal donante (no en el supremo esencial), vale al menos hasta $\Lambda=5,6$ en el escenario base, 1,8 en el de confusión por el historial y 7,3 en el de confusión no observada. Con $U$ normal, $\lambda$ no está acotado y ningún $\Lambda$ finito cumple el modelo MSM en sentido estricto, así que en ese escenario la Proposición~\ref{prop:sensibilidad} no garantiza nada; $\Lambda=2,26$ contiene $\theta$ por construcción, y lo usamos solo como escala.

\begin{table}[ht]
\centering\small
\caption{Cobertura del intervalo combinado de sensibilidad.}\label{tab:sensibilidad}
\begin{tabular}{@{}lccccc@{}}
\toprule
Escenario & $\theta/\theta_0$ & $\Lambda=1,25$ & $\Lambda=1,5$ & $\Lambda=2$ & $\Lambda=2,5$ \\
\midrule
$C$ confunde, sin $C$ & 1,30 & 87\,\% & 100\,\% & 100\,\% & 100\,\% \\
Confusión no observada & 2,26 & 0\,\% & 3\,\% & 71\,\% & 99\,\% \\
\bottomrule
\end{tabular}
\end{table}

Para calibrar $\Lambda$, en el escenario de confusión no observada $U$ entra en el canal con coeficiente 1 y la covariable observada más fuerte, $V_1$, con 0,5. Predijimos, antes de calcularlo, que el método del Lema~\ref{lem:calibracion} recuperaría el valor exacto para $U$ y que $U$ superaría a todas las covariables observadas; ambas se cumplen (tabla~\ref{tab:calibracion}). El $\Lambda$ necesario para cubrir la verdad, 2,26 (conocido aquí solo porque conocemos $\theta$), queda entre la mediana y el cuantil 90 de $V_1$: basta un confusor algo más fuerte que la covariable observada más fuerte para explicar toda la diferencia. No es el $\Lambda$ de ruptura de la sección~\ref{sec:sensibilidad}, que depende de la conclusión de interés y se calcula con los datos. En una segunda estimación con otra muestra, el modelo «grupo objetivo frente a donantes» coincide con el de canal en la mediana (2,06 frente a 1,99), pero sobrestima los extremos: 6,34 y 29,98 en los cuantiles 90 y 99, frente a 5,37 y 15,51 con el de canal y 5,38 y 14,07 exactos.

\begin{table}[ht]
\centering\small
\caption{Calibración de $\Lambda$ con el modelo de canal: cuantiles en $R$ de $\max(\lambda,1/\lambda)$.}\label{tab:calibracion}
\begin{tabular}{@{}lccc@{}}
\toprule
Variable & Mediana & Cuantil 90 & Cuantil 99 \\
\midrule
$U$, exacto & 1,99 & 5,38 & 14,07 \\
$U$, estimado & 2,03 & 5,21 & 14,14 \\
$V_1$ & 1,72 & 3,07 & 5,81 \\
$V_2$ & 1,17 & 1,53 & 2,17 \\
$V_3$ & 1,14 & 1,45 & 1,98 \\
$V_4$ & 1,12 & 1,36 & 1,70 \\
$V_5$ & 1,08 & 1,21 & 1,42 \\
\bottomrule
\end{tabular}
\end{table}

\subsection{Tamaño muestral e intervalos: Wald y bootstrap}\label{sec:res-bootstrap}

Con el estimador estándar, la desviación típica del DR es 0,912, 0,508, 0,294, 0,183 y 0,112 puntos con 3\,000, 5\,000, 10\,000, 15\,000 y 30\,000 clientes: cae como $n^{-0,91}$, mucho más deprisa que $n^{-1/2}$, y la cobertura del intervalo de Wald es del 70, 80, 88, 92 y 95,5\,\%. Las predicciones de 5\,000 y 15\,000 se registraron de antemano y se cumplieron. El cociente $\neff/n_0$ estimado sigue la misma progresión (0,02, 0,07 y 0,14 con 3\,000, 10\,000 y 30\,000, frente al 0,22 teórico). Esto sugiere un diagnóstico, aún sin contrastar: si $\neff/n_0$ sigue creciendo al aumentar la muestra, los intervalos todavía no son fiables.

Evaluamos también el bootstrap por unidad (cliente), reestimando LightGBM y el \emph{cross-fitting} en cada remuestreo. Con los pesos estándar sobrestima la dispersión entre un 35 y un 60\,\% (en 80 réplicas, un error típico bootstrap medio de 0,165 frente a una dispersión real de 0,119 con 30\,000 clientes, y de 0,446 frente a 0,322 con 10\,000), y una variante bayesiana que pondera a los clientes en lugar de copiarlos solo reduce en parte el exceso. Con los pesos de corte conocido, mucho menos ruidosos, el bootstrap bayesiano queda casi calibrado: 0,101 frente a 0,096 y 0,208 frente a 0,187 en 40 réplicas, con coberturas del 95 y del 90\,\%. Fallaron las cuatro predicciones registradas sobre el bootstrap; la evidencia sugiere que el exceso está ligado al ruido de los pesos estimados, y dejamos los detalles para la nota complementaria.
